\documentclass[11pt,a4paper]{article}

% ── Encoding & fonts ────────────────────────────────────────────────
\usepackage[utf8]{inputenc}
\usepackage[T1]{fontenc}
\usepackage{lmodern}

% ── Mathematics ─────────────────────────────────────────────────────
\usepackage{amsmath,amssymb,mathtools}

% ── Graphics & tables ──────────────────────────────────────────────
\usepackage{graphicx}
\usepackage{booktabs}

% ── Lists & formatting ─────────────────────────────────────────────
\usepackage{enumitem}

% ── Hyperlinks & metadata ──────────────────────────────────────────
\usepackage[colorlinks=true,linkcolor=blue,citecolor=blue,urlcolor=blue]{hyperref}

% ── Citations ───────────────────────────────────────────────────────
\usepackage[numbers,sort&compress]{natbib}

% ── Custom macros ───────────────────────────────────────────────────
\newcommand{\sg}{\operatorname{sg}}        % stop-gradient
\newcommand{\pos}{\operatorname{pos}}      % positional encoding
\newcommand{\mask}{\operatorname{mask}}    % mask token
\newcommand{\EMA}{\operatorname{EMA}}      % exponential moving average
\DeclareMathOperator*{\argmin}{arg\,min}

% ── Title ───────────────────────────────────────────────────────────
\title{%
  LavaJEPA: Comparing JEPA-Style Predictive Objectives\\
  on Spiking and Classical Neural Networks%
}
\author{David Strnadel}
\date{February 2026}

\begin{document}

\maketitle

\begin{abstract}
We investigate whether a JEPA-style predictive world-model objective is more
naturally realizable using spiking neural networks (neuromorphic approach
simulated in Intel Lava) than using classical neural networks on GPU.  Both
implementations solve the same self-supervised prediction task over time-series
sensory data, use the same inputs, and are evaluated with the same metrics.
This document formalizes the method, experimental setup, and evaluation
protocol.
\end{abstract}

\section{Introduction}

The Joint Embedding Predictive Architecture
(JEPA)~\cite{lecun2022path,assran2023ijepa} proposes learning
representations by predicting latent states rather than raw signal values.
This paradigm naturally involves temporal context, internal state maintenance,
and prediction---properties that closely align with the computational
principles of spiking neural networks (SNNs).

We hypothesize that the JEPA objective maps more naturally onto neuromorphic
computation, where temporal dynamics and sparse event-driven processing are
first-class citizens, than onto conventional GPU-based architectures that
process time as a tensor dimension.

To test this hypothesis, we build two minimal, parallel implementations of the
same JEPA-style predictive task: one using a classical ANN (PyTorch~\cite{paszke2019pytorch}) and one
using an SNN (Intel Lava, CPU simulation)~\cite{lava2023}.  Both are evaluated on shared
metrics covering prediction quality, latency, and computational effort.

\section{Method}\label{sec:method}

\subsection{JEPA-style predictive objective in time}

Given a time-series window $x_{1:T}$, we define two disjoint index sets:
\begin{equation}\label{eq:index-sets}
\mathcal{C} \subset \{1,\dots,T\}, \quad
\mathcal{T} \subset \{1,\dots,T\}, \quad
\mathcal{C}\cap\mathcal{T}=\emptyset,
\end{equation}
where $\mathcal{C}$ are context timesteps and $\mathcal{T}$ are target
timesteps to be predicted.

A \emph{student} (context) encoder $f_\theta$ processes only context timesteps
and produces latent representations:
\begin{equation}\label{eq:student}
z_t^{(s)} = f_\theta(x_{1:T};\,\mathcal{C})_t, \quad t\in\mathcal{C}.
\end{equation}

A \emph{teacher} (target) encoder $\bar{f}_{\bar{\theta}}$ processes the full
window:
\begin{equation}\label{eq:teacher}
z_t^{(t)} = \bar{f}_{\bar{\theta}}(x_{1:T})_t, \quad t\in\mathcal{T},
\end{equation}
with stop-gradient applied to its outputs.

A \emph{predictor} $g_\phi$ receives the student latents together with
time-position and mask tokens and predicts target latents:
\begin{equation}\label{eq:predictor}
\hat{z}_t = g_\phi\!\Bigl(\{z_u^{(s)}\}_{u\in\mathcal{C}},\;
  \pos(t),\;\mask(t)\Bigr), \quad t\in\mathcal{T}.
\end{equation}

The objective is a representation-space distance applied only to target
indices:
\begin{equation}\label{eq:loss}
\mathcal{L}_{\text{JEPA-time}}(\theta,\phi) =
\frac{1}{|\mathcal{T}|}
\sum_{t\in\mathcal{T}}
d\!\bigl(\hat{z}_t,\;\sg(z_t^{(t)})\bigr),
\end{equation}
where $d(\cdot,\cdot)$ is typically MSE or cosine distance and
$\sg(\cdot)$ denotes the stop-gradient operator.

Teacher parameters are updated by exponential moving average:
\begin{equation}\label{eq:ema}
\bar{\theta} \leftarrow \tau\,\bar{\theta} + (1-\tau)\,\theta.
\end{equation}

\subsection{Temporal masking strategies}

We evaluate three masking policies:

\begin{enumerate}[label=(\roman*)]
\item \textbf{Future-block prediction:}
  $\mathcal{C}=\{1,\dots,t_0\}$, \;
  $\mathcal{T}=\{t_0{+}1,\dots,t_0{+}k\}$.
\item \textbf{Random drop-in-window:}
  $\mathcal{T}$ is sampled uniformly within $\{1,\dots,T\}$;
  $\mathcal{C}$ is the complement.
\item \textbf{Multi-target blocks:}
  multiple disjoint target intervals inside the window.
\end{enumerate}

\section{Experimental Setup}

\subsection{ANN baseline (GPU)}

The ANN implementation follows the JEPA factorization:

\begin{itemize}[nosep]
\item \textbf{Encoder:} GRU or small Transformer producing per-timestep
  latent vectors~$z_t$.
\item \textbf{Predictor:} small MLP or Transformer conditioned on
  time-position tokens.
\item \textbf{Teacher:} EMA copy of the encoder (stop-gradient targets).
\end{itemize}

\subsection{SNN implementation (Lava CPU simulation)}

The SNN implementation mirrors the same objective using Lava processes:

\begin{itemize}[nosep]
\item \textbf{Spiking encoder:} recurrent LIF neuron population; internal
  state is stored as membrane potentials and filtered spike traces.
\item \textbf{Predictor:} spiking or hybrid module producing a continuous
  latent estimate from spike traces.
\item \textbf{Teacher:} EMA-updated copy of encoder weights (updated between
  runs).
\end{itemize}

Execution proceeds in discrete timesteps using Lava \texttt{RunSteps}.
Neuronal state is stored in \texttt{Vars} and communication occurs
via~\texttt{Ports}.

\subsection{Datasets (selection criteria)}

We consider temporal sensory datasets where prediction is natural and
measurable:

\begin{itemize}[nosep]
\item short video sequences (latent prediction of future frame
  representations),
\item audio windows (prediction of future spectral features),
\item IMU sensor streams (prediction of future motion state).
\end{itemize}

These datasets share:
\begin{itemize}[nosep]
\item continuous time structure,
\item natural next-state predictability,
\item suitability for streaming, batch-size-one evaluation.
\end{itemize}

\subsection{Evaluation metrics}

\paragraph{Prediction loss.}
The JEPA-time loss~\eqref{eq:loss} averaged over target timesteps.

\paragraph{Latency.}
\begin{itemize}[nosep]
\item Algorithmic latency: number of simulation steps per prediction.
\item Wall-clock latency: milliseconds per prediction (GPU vs.\ Lava CPU
  runtime).
\end{itemize}

\paragraph{Event metrics (SNN only).}
\begin{itemize}[nosep]
\item Total number of spikes,
\item spikes per timestep,
\item synaptic/message events (if available from Lava instrumentation).
\end{itemize}

\paragraph{Energy proxy.}
\begin{equation}\label{eq:energy}
\hat{E} =
  \alpha\, N_{\text{syn}}
  + \beta\, N_{\text{spikes}}
  + \gamma\, N_{\text{steps}},
\end{equation}
with coefficients taken from published Loihi energy measurements where
available.

\begin{table}[htbp]
\centering
\caption{Mapping JEPA concepts to ANN and SNN implementations.}
\label{tab:mapping}
\smallskip
\begin{tabular}{@{}lll@{}}
\toprule
\textbf{JEPA concept} & \textbf{ANN (PyTorch)} & \textbf{SNN (Lava)} \\
\midrule
Predict embeddings, not signal
  & Predict latent $z_t$
  & Predict latent from spike traces \\
Mask context/targets
  & Mask timesteps
  & Same masking via input gating \\
Predictor with position tokens
  & Time-position embeddings
  & Time-index input channel via Ports \\
EMA teacher + stop-gradient
  & EMA encoder copy
  & EMA copy between runs \\
Execution over time
  & Tensor time dimension
  & \texttt{RunSteps}, state in \texttt{Vars} \\
Energy/latency rationale
  & Measure wall-clock latency
  & Measure spikes, events, timesteps \\
\bottomrule
\end{tabular}
\end{table}

\section{Design Principles}

This section summarizes the key design principles derived from the JEPA
and neuromorphic computing literature that motivate the choices made in
Section~\ref{sec:method}.

\subsection{Predicting embeddings, not signal}

JEPA explicitly predicts target-encoder embeddings rather than masked raw
input data.  The masking and prediction both occur in representation space.
This is critical because the targets are semantic latent features rather
than low-level signal values.  V-JEPA~\cite{bardes2024vjepa} extends this
to video, and VL-JEPA~\cite{zhai2025vljepa}
to text embeddings, reinforcing the principle that embedding-space
supervision is preferred over signal reconstruction.

\subsection{Context/target split and temporal translation}

In I-JEPA~\cite{assran2023ijepa}, multiple target blocks are sampled in representation space while
a context block is chosen to avoid overlap with targets.  In
V-JEPA~\cite{bardes2024vjepa},
masked patches correspond to dropped spatio-temporal regions.

For our temporal setting, two-dimensional spatial blocks translate
naturally into disjoint time-index sets $\mathcal{C}$ and $\mathcal{T}$
(Section~\ref{sec:method}), and two-dimensional positional encodings are
replaced with scalar time-position tokens.

\subsection{Collapse prevention}

The JEPA literature provides two explicit anti-collapse mechanisms:

\begin{enumerate}[nosep,label=(\roman*)]
\item EMA teacher encoder with stop-gradient
  (I-JEPA~\cite{assran2023ijepa}, V-JEPA~\cite{bardes2024vjepa}),
\item batch-level regularization such as InfoNCE alignment and uniformity
  (VL-JEPA~\cite{zhai2025vljepa}).
\end{enumerate}

For a clean comparison between ANN and SNN we adopt approach~(i), as it
requires no batch-level contrastive machinery and is compatible with
streaming, batch-size-one evaluation.

\subsection{Minimal architecture requirement}

The JEPA papers do not require Vision Transformers specifically.  The
architectural requirement reduces to the three-component factorization
(encoder, predictor, EMA teacher) described in
Section~\ref{sec:method}.  For time series this can be instantiated with
a small GRU/LSTM or Transformer encoder, a small MLP predictor, and an
EMA copy of the encoder.

\subsection{Neuromorphic workload constraints}

The Loihi literature~\cite{davies2018loihi,orchard2021loihi2} highlights
that conventional feedforward networks
benefit little from neuromorphic hardware.  Advantages appear when:

\begin{itemize}[nosep]
\item computation is sparse and event-driven,
\item recurrence and temporal state are central,
\item batch-size-one, low-latency streaming inference is meaningful.
\end{itemize}

The JEPA predictive objective over temporal data satisfies all three
criteria, making it a suitable workload for neuromorphic evaluation.

\subsection{Event counts as energy proxies}

Loihi energy consumption scales with the number of virtual timesteps and
decreases with input sparsity.  Valid hardware-proxy metrics are therefore
the number of simulation steps, total spikes, and total synaptic events.
These are formalized in the energy proxy~\eqref{eq:energy}.

\subsection{SNN architectural implications from Lava}

Lava defines computation as \texttt{Processes} with internal state
(\texttt{Vars}) communicating via \texttt{Ports}.  Execution proceeds in
discrete timesteps via \texttt{RunSteps}.  The SNN implementation is thus
structured as:

\begin{itemize}[nosep]
\item encoder process (recurrent LIF population),
\item predictor process,
\item output readout process,
\item explicit timestep-based execution loop.
\end{itemize}

\subsection{Predictive coding and local learning rules}

Spiking predictive-coding literature~\cite{rao1999predictive} shows that
prediction errors can be
represented by explicit error neurons, membrane-potential differences, or
implicitly in network dynamics.  LIF neuron dynamics provide natural
temporal state through membrane potentials:
\begin{equation}\label{eq:lif}
\tau \frac{dV}{dt} = -V + I(t),
\end{equation}
with spike traces updated by exponential decay and spike increments.
Predictive-coding learning rules offer local-update alternatives to full
BPTT, which is important because training recurrent SNNs via BPTT is
computationally expensive.

% ====================================================================
\section{Implementation Specification}\label{sec:impl-spec}
% ====================================================================

This section serves as a technical specification (``RFC'') for both
implementations.  It defines invariants, interfaces, baseline
configurations, and measurement procedures so that the comparison remains
fair, reproducible, and free from accidental degrees of freedom.

% --------------------------------------------------------------------
\subsection{Invariants and task contract}\label{sec:invariants}
% --------------------------------------------------------------------

The following invariants \textbf{must not} differ between the ANN and
SNN implementations:

\begin{enumerate}[label=\textbf{I\arabic*.},nosep]
\item \textbf{Same task.}
  Both models receive the same time-series window $x_{1:T}$, the same
  context/target split $(\mathcal{C},\mathcal{T})$, and are evaluated
  with the same loss over~$\mathcal{T}$.
\item \textbf{Same latent target.}
  Predictions $\hat{z}_t$ are compared against teacher-encoder outputs
  $z_t^{(t)}$ produced by conceptually the same EMA teacher
  (Eq.~\ref{eq:ema}).
\item \textbf{Same inference regime.}
  All evaluation metrics are collected at batch size $B=1$ in a
  streaming (or fixed-window) mode, identically for both sides.
\item \textbf{Same latency definition.}
  Latency is reported as (a)~wall-clock milliseconds per one target-block
  prediction and (b)~number of simulation steps.
\end{enumerate}

\paragraph{Task contract (API).}
Both implementations must expose the following interface:

\begin{center}
\begin{tabular}{@{}lll@{}}
\toprule
\textbf{Name} & \textbf{Shape} & \textbf{Description} \\
\midrule
Input $x$
  & $B \times T \times D$
  & Raw time-series window \\
Mask $m$
  & $B \times T$
  & Binary; $m_t=1$ iff $t\in\mathcal{T}$ \\
Prediction $\hat{z}$
  & $B \times T \times H$
  & Filled only for target $t$ \\
Teacher target $z^{(t)}$
  & $B \times T \times H$
  & Evaluated only for target $t$ \\
Loss $\mathcal{L}$
  & scalar
  & $\frac{1}{|\mathcal{T}|}\sum_{t\in\mathcal{T}}
     \lVert \hat{z}_t - \sg(z_t^{(t)}) \rVert^2$ \\
\bottomrule
\end{tabular}
\end{center}

This contract is the ground truth; architectures may change, but the
contract must not.

% --------------------------------------------------------------------
\subsection{Baseline configurations}\label{sec:baselines}
% --------------------------------------------------------------------

\paragraph{Comparable-capacity rule.}
To ensure a fair comparison, both implementations fix the latent
dimension~$H$ and approximately match the total number of trainable
parameters.  Both~$H$ and parameter counts are reported explicitly.

\subsubsection{ANN baseline}

\begin{itemize}[nosep]
\item \textbf{Encoder:} 1-layer GRU, hidden size $H=128$.
\item \textbf{Predictor:} 2-layer MLP ($H\!\to\!H\!\to\!H$) with
  time-position embedding concatenated at the input.
\item \textbf{Teacher:} EMA copy of encoder, $\tau=0.996$ (fixed).
\item \textbf{Loss:} MSE.
\item \textbf{Context length:} $T$ timesteps (fixed for all runs).
\end{itemize}

\subsubsection{SNN baseline}

\begin{itemize}[nosep]
\item \textbf{Encoder:} Recurrent LIF population, $H=128$ neurons,
  with surrogate-gradient training in PyTorch (see
  Section~\ref{sec:training}).
\item \textbf{Spike trace:} Exponential low-pass filter of output spikes:
  \begin{equation}\label{eq:spike-trace}
    r_t = \alpha\, r_{t-1} + (1-\alpha)\, s_t,
  \end{equation}
  where $s_t\in\{0,1\}^H$ is the spike vector and $\alpha$ is the
  trace decay constant.
\item \textbf{Readout:} Linear projection from spike trace to latent
  space:
  \begin{equation}\label{eq:readout}
    \hat{z}_t = W_{\text{ro}}\, r_t + b_{\text{ro}}.
  \end{equation}
\item \textbf{Predictor:} 2-layer spiking or hybrid MLP (same topology
  as ANN predictor).
\item \textbf{Teacher:} EMA copy of encoder weights, $\tau=0.996$.
\item \textbf{Loss:} MSE (computed on readout $\hat{z}_t$).
\item \textbf{Context length:} same $T$ as ANN.
\end{itemize}

\paragraph{Parameter reporting.}
For each configuration we report:
\begin{itemize}[nosep]
\item total trainable parameters (\#params),
\item FLOPs per forward pass (ANN),
\item spike events, synaptic events, and simulation steps (SNN).
\end{itemize}

% --------------------------------------------------------------------
\subsection{Latent space definition}\label{sec:latent-def}
% --------------------------------------------------------------------

A frequent source of hidden divergence is ambiguity about \emph{where} in
the network the latent~$z_t$ is taken and how its scale is controlled.
This subsection pins down those choices.

\paragraph{Latent tap point.}
In both ANN and SNN, the ``latent'' $z_t$ is the \textbf{direct output of
the encoder}, with no additional projection head.  Concretely:
\begin{itemize}[nosep]
\item \textbf{ANN:} $z_t = h_t \in \mathbb{R}^H$, where $h_t$ is the
  GRU hidden state at step~$t$.
\item \textbf{SNN:} $z_t = \hat{z}_t = W_{\text{ro}}\,r_t +
  b_{\text{ro}}$ (Eq.~\ref{eq:readout}), the linear readout of the
  spike trace.
\end{itemize}
Adding a projection head is explicitly \textbf{not} done; if future work
adds one, it must be identical in topology and dimension for both sides.

\paragraph{Latent normalization.}
Teacher targets $z_t^{(t)}$ are \emph{not} L2-normalized before the loss
computation (the loss is raw MSE).  However, we log the per-batch mean
embedding norm $\mathbb{E}[\lVert z_t \rVert_2]$ for both student and
teacher to detect norm drift or collapse.

\paragraph{Readout compatibility.}
The SNN readout (Eq.~\ref{eq:readout}) and the ANN GRU output both
produce vectors in $\mathbb{R}^H$.  The readout weight matrix
$W_{\text{ro}} \in \mathbb{R}^{H\times H}$ and bias
$b_{\text{ro}} \in \mathbb{R}^H$ are trained jointly with the SNN
encoder via surrogate gradients and exported to Lava as a
\texttt{Dense} layer after the spike-trace computation.

% --------------------------------------------------------------------
\subsection{Tokenization and masking}\label{sec:tokens}
% --------------------------------------------------------------------

Positional and mask tokens must be defined identically in both
implementations to prevent divergence.

\paragraph{Time-position token $\pos(t)$.}
Sinusoidal encoding of the timestep index, producing a vector
$e_t\in\mathbb{R}^H$:
\begin{equation}\label{eq:pos-enc}
e_{t,2k}   = \sin\!\bigl(t / 10000^{2k/H}\bigr), \quad
e_{t,2k+1} = \cos\!\bigl(t / 10000^{2k/H}\bigr).
\end{equation}
The same $e_t$ values are used by both ANN and SNN predictors.

\paragraph{Mask token $\mask(t)$.}
\begin{itemize}[nosep]
\item \textbf{ANN:} a learned embedding vector $m_{\text{tok}} \in
  \mathbb{R}^H$ replaces student latents at target positions.
\item \textbf{SNN:} an extra binary input channel implemented as a
  \emph{deterministic constant bias current}: $I_{\text{mask}}(t) =
  c_{\text{mask}}$ for $t\in\mathcal{T}$ and $0$ otherwise, injected
  via a dedicated Lava \texttt{Dense} connection with a single learned
  weight column.  No Poisson noise is added; the signal is identical
  across runs for a given mask.  The bias magnitude $c_{\text{mask}}$
  is a fixed hyperparameter (default: $c_{\text{mask}} = 1.0$,
  tunable only during development, frozen before final runs).
\end{itemize}

% --------------------------------------------------------------------
\subsection{Masking parameters}\label{sec:mask-params}
% --------------------------------------------------------------------

All three masking policies from Section~\ref{sec:method} share a common
window length~$T$.  The \emph{baseline} policy used for the primary
comparison is future-block prediction; the other two are ablations.

\paragraph{(i) Future-block prediction (baseline).}
Context comprises the first 75\,\% of the window, target the remaining
25\,\%:
\[
  t_0 = \lfloor 0.75\,T \rfloor, \quad
  k   = T - t_0, \quad
  \mathcal{C}=\{1,\dots,t_0\}, \quad
  \mathcal{T}=\{t_0{+}1,\dots,T\}.
\]
For $T=128$: $t_0=96$, $k=32$.

\paragraph{(ii) Random drop-in-window.}
Each timestep is independently assigned to~$\mathcal{T}$ with probability
$p=0.25$ (without replacement constraint; sampling uses the training-order
seed).  $\mathcal{C}$ is the complement.

\paragraph{(iii) Multi-target blocks.}
Two contiguous target blocks of length $\lfloor 0.125\,T \rfloor = 16$
each, placed uniformly at random with a minimum gap of $\lfloor 0.0625\,T
\rfloor = 8$ timesteps between them.  $\mathcal{C}$ is the complement.

All mask indices are generated deterministically from the per-window seed
so that ANN and SNN receive identical masks.

% --------------------------------------------------------------------
\subsection{Lava process graph}\label{sec:lava-graph}
% --------------------------------------------------------------------

The SNN inference network in Lava is a fixed directed graph of
\texttt{Process} instances connected via \texttt{Ports}.
Figure~\ref{fig:lava-graph} defines the topology; the table below
specifies each process, its I/O port shapes, data types, and the Lava
class used.

\begin{figure}[htbp]
\centering
\small
\begin{verbatim}
  InputSource ──s_out──▶ DenseIn ──a_out──▶ LIF_enc ──s_out──┐
                                                              │
  MaskSource ──s_out──▶ DenseMask ──a_out──────────────────▶(+)
                                                              │
  PosSource  ──s_out──▶ DensePos ──a_out───────────────────▶(+)
                                                              │
                                              ◀──s_in── DensePred
                                                          │
                                                       a_out
                                                          │
                                                       LIF_pred
                                                          │
                                                       s_out
                                                          │
                                                    DenseReadout
                                                          │
                                                       a_out
                                                          │
                                                    OutputSink
\end{verbatim}
\caption{Lava process graph for SNN inference.  Arrows indicate port
connections; ``\texttt{(+)}'' denotes current summation at the
\texttt{a\_in} port of \texttt{LIF\_pred}.  Three input channels feed the
predictor: encoder spikes, mask signal, and positional signal.}
\label{fig:lava-graph}
\end{figure}

\begin{center}
\renewcommand{\arraystretch}{1.12}
\begin{tabular}{@{}llll@{}}
\toprule
\textbf{Process} & \textbf{Lava class} & \textbf{Shape / Port} &
  \textbf{Notes} \\
\midrule
\texttt{InputSource}
  & \texttt{RingBuffer}
  & out: $(D, T)$
  & Pre-loaded window $x_{1:T}$ \\
\texttt{MaskSource}
  & \texttt{RingBuffer}
  & out: $(1, T)$
  & $c_{\text{mask}}$ for $t\!\in\!\mathcal{T}$, else 0 \\
\texttt{PosSource}
  & \texttt{RingBuffer}
  & out: $(H, T)$
  & $\pos(t)$ (Eq.~\ref{eq:pos-enc}), all $t$ \\
\texttt{DenseIn}
  & \texttt{Dense}
  & $W \!\in\! \mathbb{R}^{H\times D}$
  & Input projection \\
\texttt{DenseMask}
  & \texttt{Dense}
  & $w \!\in\! \mathbb{R}^{H\times 1}$
  & Mask channel weight \\
\texttt{DensePos}
  & \texttt{Dense}
  & $W \!=\! I_{H}$ (identity)
  & Positional channel (pass-through) \\
\texttt{LIF\_enc}
  & \texttt{LIF}
  & $(H,)$; $\texttt{dv}\!=\!1{-}\beta$
  & Encoder neurons \\
\texttt{DensePred}
  & \texttt{Dense}
  & $W \!\in\! \mathbb{R}^{H\times H}$
  & Predictor layer~1 \\
\texttt{LIF\_pred}
  & \texttt{LIF}
  & $(H,)$
  & Predictor neurons \\
\texttt{DenseReadout}
  & \texttt{Dense}
  & $W_{\text{ro}} \!\in\! \mathbb{R}^{H\times H}$
  & Readout (Eq.~\ref{eq:readout}) \\
\texttt{OutputSink}
  & \texttt{RingBuffer}
  & in: $(H, T)$
  & Collects $\hat{z}_t$ \\
\bottomrule
\end{tabular}
\end{center}

\paragraph{Predictor input channels.}\label{sec:pred-channels}
The SNN predictor (\texttt{LIF\_pred}) receives exactly three additive
current contributions at its \texttt{a\_in} port:
\begin{enumerate}[nosep,label=(\alph*)]
\item \textbf{Encoder spikes.}
  \texttt{LIF\_enc.s\_out} $\to$ \texttt{DensePred} $\to$
  \texttt{LIF\_pred.a\_in}.  This is the primary signal carrying the
  student's latent representation.
\item \textbf{Positional signal.}
  \texttt{PosSource} emits the sinusoidal vector
  $\pos(t) \in \mathbb{R}^H$ (Eq.~\ref{eq:pos-enc}) as a graded
  (non-spiking) current at every timestep.  \texttt{DensePos} is an
  identity-weight \texttt{Dense} process ($W = I_H$) that relays
  $\pos(t)$ to \texttt{LIF\_pred.a\_in}.  The numerical values are
  \emph{identical} to those added to the ANN predictor input.
\item \textbf{Mask signal.}
  \texttt{MaskSource} emits a scalar $c_{\mask}$ at masked positions
  ($t \in \mathcal{T}$) and 0 elsewhere.  \texttt{DenseMask} ($w \in
  \mathbb{R}^{H \times 1}$) broadcasts this scalar to all $H$
  predictor neurons.  This is the SNN analogue of the ANN's learned
  \texttt{[MASK]} embedding.
\end{enumerate}

\noindent\textbf{Exclusivity rule.}  No other input signals may be
connected to \texttt{LIF\_pred.a\_in} beyond (a)--(c).  Any additional
``helper'' currents (e.g.\ learnable bias, noise injection, activity
regularization signals) that have no ANN counterpart constitute a
deviation and must be treated as a separate ablation experiment
(Section~\ref{sec:equivalence}).

\paragraph{Spike/event collection.}
A Lava \texttt{Monitor} is attached to \texttt{LIF\_enc.s\_out} and
\texttt{LIF\_pred.s\_out} to record per-timestep spike vectors.  Spike
count is $\sum_t \sum_i s_{i,t}$ over both populations.  Synaptic events
are computed post-hoc as $N_{\text{syn}} = \sum_t \lVert s_t \rVert_0
\cdot N_{\text{fan-out}}$ for each \texttt{Dense} layer, where
$N_{\text{fan-out}}$ is the number of non-zero entries in the
corresponding weight column.  If Lava provides native event counters in
future versions, those are used instead.

\paragraph{Execution.}
The graph is executed as a single \texttt{RunSteps(num\_steps=T)} call
using \texttt{Loihi2SimCfg(select\_tag="floating\_pt")}.  After the run,
data is retrieved from sinks and monitors, then \texttt{stop()} is
called.

% --------------------------------------------------------------------
\subsection{Training protocol}\label{sec:training}
% --------------------------------------------------------------------

\paragraph{Separation of training and measurement.}
Training and inference/measurement are performed in separate frameworks:

\begin{center}
\begin{tabular}{@{}lll@{}}
\toprule
& \textbf{Training} & \textbf{Inference / Measurement} \\
\midrule
ANN & PyTorch (GPU/CPU) & PyTorch (CPU, no grad) \\
SNN & PyTorch + surrogate gradients & Intel Lava (CPU simulation) \\
\bottomrule
\end{tabular}
\end{center}

SNN weights are trained in PyTorch using surrogate-gradient
back-propagation through time, then exported to Lava
\texttt{Dense}/\texttt{LIF} processes for measurement.  This avoids
low-level training issues in Lava while keeping the inference
measurement authentic.

\paragraph{SNN neuron model (training).}\label{sec:snn-neuron}
During training, LIF neurons are simulated in discrete time with
Euler integration ($\Delta t = 1$ step):
\begin{align}\label{eq:lif-discrete}
  V_i[t] &= \beta\, V_i[t{-}1]\,(1-s_i[t{-}1]) + I_i[t], \\
  s_i[t]  &= \Theta(V_i[t] - v_{\text{th}}), \notag
\end{align}
where $\beta = \exp(-\Delta t / \tau_{\text{mem}})$ is the membrane
decay, $\Theta$ is the Heaviside step function (replaced by the
surrogate during back-propagation), and the reset is multiplicative
(hard reset to zero on spike).

\emph{Fixed parameters:} $\tau_{\text{mem}} = 10$ steps,
$v_{\text{th}} = 1.0$, no synaptic time constant
($\tau_{\text{syn}} = 0$; current-based synapses).

\emph{Surrogate gradient:} Fast
sigmoid~\cite{neftci2019surrogate,zenke2021remarkable}
\[
  \tilde{\Theta}'(x) = \frac{1}{(1 + k|x|)^2}, \quad k = 25.
\]
No other surrogate is used; $k$ is fixed before final runs.

\paragraph{Weight export to Lava.}
Trained \texttt{torch.nn.Linear} weights are mapped one-to-one to
Lava \texttt{Dense} weight matrices ($W_{\text{Lava}} = W_{\text{torch}}^{\!\top}$,
transposed to match Lava's $(N_{\text{out}},N_{\text{in}})$
convention).  LIF parameters are set as:
$\texttt{du} = 0$, $\texttt{dv} = 1 - \beta$,
$\texttt{vth} = v_{\text{th}}$.  A round-trip unit test verifies that
Lava and PyTorch produce identical spike trains (tolerance:
$\leq 1$ spike mismatch per 1\,000 steps per neuron) before any
measurement run.

\paragraph{Training schedule.}
\begin{itemize}[nosep]
\item Fixed number of training steps $N_{\text{train}}$ (identical for
  both).
\item Learning rate: cosine schedule with linear warm-up (5\% of
  $N_{\text{train}}$).
\item Optimizer: AdamW, $\text{lr}_{\max}$, $\beta=(0.9,\,0.999)$,
  weight decay $10^{-4}$.
\item EMA momentum $\tau=0.996$ (fixed, not annealed).
\end{itemize}

\paragraph{Teacher target handling.}\label{sec:teacher-handling}
The default configuration \emph{pre-computes} all teacher targets
$z_t^{(t)}$ for $t \in \mathcal{T}$ across the entire test set before
any measurement begins.  Pre-computed targets are stored as
\texttt{float32} \texttt{.npz} files and verified by SHA-256 hash.
During the timed latency measurement, the teacher network is
\emph{not executed}; targets are loaded from disk, and latency
component~(2) (Section~\ref{sec:measurement}) is reported as zero.

An optional \emph{online-teacher variant} re-runs the teacher at
measurement time.  Rules for this variant:
\begin{itemize}[nosep]
\item The teacher model must be identical for ANN and SNN (same weights,
  same architecture, same framework---PyTorch).
\item Teacher forward latency is recorded separately in component~(2)
  and is \emph{not} added to the primary encoder+predictor comparison
  in component~(1).
\item This variant receives its own experiment identifier and result
  rows.
\end{itemize}

\paragraph{Dataset.}\label{sec:dataset}
The primary dataset is the \textbf{UCI Human Activity Recognition Using
Smartphones} dataset~\cite{anguita2013har} (Anguita et~al., 2013).  Raw 3-axis accelerometer and
3-axis gyroscope signals ($D=6$) are sampled at 50\,Hz.  Each window spans
$T_{\text{sec}}=2.56$\,s, yielding $T=128$ samples per window.

\emph{Preprocessing.}
Per-channel z-score normalization ($\mu=0$, $\sigma=1$) is computed on the
training set only and applied identically to validation and test sets.

\emph{Data split.}
The split is \textbf{by subject ID} (not by random windows) to prevent
intra-subject leakage: 70\,\% train / 15\,\% validation / 15\,\% test
subjects (21 / 4 / 5 subjects out of 30).
This yields $N_{\text{tr}} = 7209$ training windows,
$N_{\text{val}} = 1373$ validation windows, and
$N_{\text{te}} = 1716$ test windows
(10\,299 total).
The split seed is logged (see Section~\ref{sec:repro}).

\emph{Secondary modality.}
Audio feature streams (e.g.\ LibriSpeech 80-dim log-mel, 16\,kHz,
$T=128$ frames at 10\,ms hop) serve as a secondary modality for
generalization checks.  Video is reserved for an optional extension.

\paragraph{Data pipeline implementation.}\label{sec:data-pipeline}
The following steps form the canonical data pipeline; any deviation
must be documented and justified.

\begin{enumerate}[nosep,label=\textbf{D\arabic*.}]
\item \textbf{Download.}
  The dataset is obtained via the official UCI ML Repository URL and
  verified by SHA-256 hash of the archive.  The download script is
  version-controlled in \texttt{data/download.py}.
\item \textbf{Subject split.}
  Subject IDs are sorted, then deterministically partitioned using the
  data-split seed (R1).  The mapping
  \texttt{subject\_id $\to$ \{train, val, test\}} is saved as a JSON
  file and logged with each run.
\item \textbf{Windowing.}
  Fixed-length windows of $T=128$ samples are extracted with no overlap
  (stride $=T$).  Partial windows at end of recording are discarded.
\item \textbf{Normalization.}
  Per-channel mean $\mu_d$ and standard deviation $\sigma_d$ are computed
  \emph{on training windows only}, then applied to all splits:
  $\tilde{x}_{t,d} = (x_{t,d} - \mu_d)/\sigma_d$.  The statistics
  $(\mu, \sigma)$ are saved to disk and loaded (not recomputed) during
  validation/test.
\item \textbf{Mask generation.}
  For each window, the per-window seed is derived as
  $\texttt{seed}_w = \texttt{hash}(\texttt{split\_seed},\,
  \texttt{window\_index})$.  This seed deterministically generates
  $(\mathcal{C}, \mathcal{T})$ according to the active masking policy
  (Section~\ref{sec:mask-params}).  A unit test asserts that calling the
  mask generator twice with the same seed produces bit-identical output.
\item \textbf{Iterator.}
  A single \texttt{WindowDataset} class (shared by both ANN and SNN)
  yields tuples $(x, m, \texttt{seed}_w)$ in a fixed, reproducible
  order.  During measurement, iteration is sequential (no shuffling).
\end{enumerate}

% --------------------------------------------------------------------
\subsection{Experimental design and statistical plan}%
\label{sec:exp-design}
% --------------------------------------------------------------------

\paragraph{Primary hypothesis.}
\emph{An SNN (LIF neurons in Lava) achieves lower or comparable
JEPA-time prediction loss
$\mathcal{L}$ (Eq.~\ref{eq:loss}) to a parameter-matched ANN
(GRU), at lower latency, on streaming time-series data.}

\noindent The \textbf{primary endpoint} is the paired difference
$\Delta\mathcal{L}_s = \mathcal{L}_s^{\text{SNN}}
 - \mathcal{L}_s^{\text{ANN}}$
computed over the full test set for each seed~$s$, holding masking
policy, data split, and initialization strategy fixed.

\paragraph{Design matrix.}
The study has two tiers:

\begin{center}
\renewcommand{\arraystretch}{1.12}
\begin{tabular}{@{}llccccc@{}}
\toprule
\textbf{Tier} & \textbf{Purpose} & \textbf{Seeds}
  & \textbf{Policies} & \textbf{Ablations} & \textbf{Substrates} & \textbf{Runs} \\
\midrule
Primary
  & Hypothesis test
  & $S=10$ & $P=3$ & --- & 2 & 60 \\
Mechanism ablation
  & Artefact control
  & $S=10$ & baseline only & $A=3$ & 2 & 60 \\
\midrule
\multicolumn{6}{@{}l}{\textbf{Total}}
  & \textbf{120} \\
\bottomrule
\end{tabular}
\end{center}

\noindent\emph{Primary grid} ($10 \times 3 \times 2 = 60$ runs):
3~masking policies (future-block, random-drop, multi-block),
each with 10~independent seeds (init + train-order;
data-split seed fixed), both substrates.

\noindent\emph{Mechanism ablations} ($S{=}10 \times A{=}3
\text{ ablation conditions} \times 2 \text{ substrates} = 60$ runs):
baseline masking only (future-block), 3~ablation conditions:
\begin{enumerate}[nosep,label=\textbf{A\arabic*.}]
\item \textbf{NoPos:} positional encoding $\pos(t)$ zeroed in both
  branches.  Tests whether the benefit is driven by temporal position
  information rather than the substrate.
\item \textbf{NoMaskTok:} mask token disabled (ANN: mask embedding $= 0$;
  SNN: $c_{\mask} = 0$).  Tests whether the mask signal leaks target
  information or is essential for convergence.
\item \textbf{OnlineTeacher:} teacher forward computed online at
  measurement time (identical PyTorch model for both substrates).
  Tests whether pre-computed targets introduce artefacts.
\end{enumerate}

\paragraph{Seed structure.}
Each seed~$s \in \{1, \dots, 10\}$ determines:
(a)~parameter initialization (\texttt{torch.manual\_seed(s)}),
(b)~training-order shuffle.
The data-split seed is held \emph{fixed} across all seeds
(same subjects in train/val/test) to enable paired analysis.
This yields 10~paired observations per comparison.

\paragraph{Statistical analysis plan (pre-registered).}
\label{par:stat-plan}
All analyses are specified \emph{before any data collection}.

\begin{enumerate}[nosep,label=\textbf{S\arabic*.}]
\item \textbf{Primary comparison.}
  For each masking policy~$p$, compute
  $\Delta\mathcal{L}_{s,p}
   = \mathcal{L}_{s,p}^{\text{SNN}} - \mathcal{L}_{s,p}^{\text{ANN}}$
  for $s = 1,\dots,10$.
  Report: mean, std, \textbf{95\,\% confidence interval} of
  $\Delta\mathcal{L}$ (paired $t$-interval), and
  \textbf{Cohen's~$d$} for paired differences.
\item \textbf{Primary test.}
  Two-sided paired $t$-test on $\{\Delta\mathcal{L}_{s,p}\}_{s=1}^{10}$
  for each policy~$p$.
  Significance level $\alpha = 0.05$.
\item \textbf{Robustness check.}
  Wilcoxon signed-rank test on the same paired differences
  (non-parametric; does not assume normality of $\Delta\mathcal{L}$).
\item \textbf{Multiple-comparison correction.}
  Holm--Bonferroni correction across all $P + A = 6$ hypothesis tests
  (3~primary policies + 3~ablation conditions).
  Adjusted $p$-values and family-wise error rate (FWER $\leq 0.05$)
  are reported.
\item \textbf{Effect-size reporting.}
  Cohen's~$d = \bar{\Delta} / s_{\Delta}$ with 95\,\% CI
  (bootstrap, 10\,000~resamples).
  Interpretation: $|d| < 0.2$ negligible, $0.2$--$0.5$ small,
  $0.5$--$0.8$ medium, $> 0.8$ large.
\item \textbf{Latency comparison.}
  Paired median wall-clock latency $\tilde{t}^{\text{SNN}}$ vs.\
  $\tilde{t}^{\text{ANN}}$ with IQR; Wilcoxon test across seeds.
\end{enumerate}

\noindent All analyses are implemented in a single
\texttt{analysis/stats.py} script, version-controlled at
\texttt{freeze-v1} (Section~\ref{sec:config}).

% --------------------------------------------------------------------
\subsection{Measurement protocol}\label{sec:measurement}
% --------------------------------------------------------------------

All metrics are evaluated on the held-out test set at batch size $B=1$
in streaming mode.

\paragraph{Metric checklist.}
\nopagebreak
\begin{center}
\renewcommand{\arraystretch}{1.15}
\begin{tabular}{@{}p{3.2cm}p{4.8cm}p{4.2cm}@{}}
\toprule
\textbf{Metric} & \textbf{Definition} & \textbf{Procedure} \\
\midrule
Prediction loss
  & $\frac{1}{|\mathcal{T}|}\sum_{t\in\mathcal{T}}
     \lVert\hat{z}_t - \sg(z_t^{(t)})\rVert^2$
  & Average over test windows \\
Wall-clock latency
  & ms per target-block prediction
  & 100 warm-up calls, then median of 1\,000 timed calls; CPU pinned \\
Algorithmic latency
  & \#forward steps (ANN) /
    \#\texttt{RunSteps} (SNN)
  & Logged per window \\
Spike count
  & $\sum_t \sum_i s_{i,t}$
  & SNN only; summed over all layers \\
Synaptic events
  & \#non-zero weight look-ups triggered by spikes
  & If Lava probe available; otherwise marked ``N/A'' \\
Energy proxy $\hat{E}$
  & Eq.~\eqref{eq:energy}
  & Coefficients from published Loihi measurements;
    if unavailable, omit $\hat{E}$ and report raw counts only \\
\#Parameters
  & Total trainable weights
  & Reported for both \\
Latent variance
  & $\operatorname{Var}_{t}[\hat{z}_t]$ over test set
  & Collapse detector (Section~\ref{sec:failure}) \\
\bottomrule
\end{tabular}
\end{center}

\paragraph{Measurement fairness policy.}\label{sec:fairness}
Wall-clock latency is hardware-dependent; comparing ANN on GPU against
SNN on CPU conflates algorithmic efficiency with platform capability.
We therefore report \textbf{two} latency configurations:

\begin{enumerate}[label=(\alph*),nosep]
\item \textbf{Platform-neutral (primary):} ANN inference on CPU
  (single-thread, \texttt{torch.set\_num\_threads(1)}) vs.\ SNN
  inference on CPU (Lava \texttt{floating\_pt}).  Both run on the same
  machine, same core pinning.  This is the primary wall-clock
  comparison.
\item \textbf{Practical throughput (supplementary):} ANN inference on
  GPU (if available) reported separately as a throughput reference.
  This number is \emph{not} directly compared against SNN-CPU but
  contextualizes real-world deployment trade-offs.
\end{enumerate}

\noindent All latency numbers specify: hardware (CPU model, GPU model),
OS, thread count, warm-up iterations, and number of timed repetitions.

\paragraph{Latency decomposition.}
Wall-clock latency is further broken down into three components,
reported separately:
\begin{enumerate}[nosep,label=(\arabic*)]
\item \textbf{Encoder + predictor forward:} the core computation
  (GRU forward / Lava \texttt{RunSteps}).
\item \textbf{Teacher forward:} teacher-encoder inference on the full
  window (needed only if teacher targets are recomputed at test time;
  if pre-computed, this is zero and noted as such).
\item \textbf{Runtime overhead:} data loading into process buffers
  (Lava \texttt{RingBuffer} setup), monitor retrieval, Python-level
  orchestration.  Measured as total minus component~(1).
\end{enumerate}

\paragraph{No-hidden-compute rule.}
The timed section (component~(1) above) must contain \emph{only} the
neural-network forward pass.  Specifically:
\begin{itemize}[nosep]
\item \textbf{Excluded} from timed section: process-graph allocation
  (\texttt{Lava.create()}, \texttt{nn.Module.\_\_init\_\_}),
  \texttt{RingBuffer} data loading, \texttt{Monitor} attachment,
  Python \texttt{for}-loop orchestration that schedules timesteps
  (Lava handles this internally), memory allocation for output buffers,
  and any logging / file I/O.
\item \textbf{Included} in timed section: ANN---\texttt{model.forward()};
  SNN---\texttt{run\_cfg.start()} through \texttt{run\_cfg.pause()}
  (i.e.\ a single \texttt{RunSteps(num\_steps=T)} call).
\item \textbf{Warm-up protocol:} before the timed loop, execute at least
  $W = 5$ untimed forward passes.  The warm-up must include the full
  graph setup so that JIT compilation (PyTorch) and process-model
  compilation (Lava) are completed.  Record $W$ in metadata.
\end{itemize}
Violating this rule invalidates the latency comparison.

% --------------------------------------------------------------------
\subsection{Observability points}\label{sec:observability}
% --------------------------------------------------------------------

To enable post-hoc auditing and debugging, every measurement run logs the
following quantities at the specified granularity.  All tensors are stored
as \texttt{float32} NumPy arrays in \texttt{.npz} archives, one file per
test window.

\paragraph{Per-timestep observables (logged for every $t \in \{1,\dots,T\}$).}
\begin{enumerate}[nosep,label=\textbf{O\arabic*.}]
\item $z_t^{(s)} \in \mathbb{R}^H$ --- student encoder output (ANN:
  GRU hidden; SNN: readout $W_{\text{ro}}\,r_t + b_{\text{ro}}$).
  Logged for $t \in \mathcal{C}$ only.
\item $z_t^{(t)} \in \mathbb{R}^H$ --- teacher encoder output.
  Logged for $t \in \mathcal{T}$ only.
\item $\hat{z}_t \in \mathbb{R}^H$ --- predictor output.
  Logged for $t \in \mathcal{T}$ only.
\item $s_t \in \{0,1\}^H$ --- spike vector (SNN only), both encoder and
  predictor populations.
\item $r_t \in \mathbb{R}^H$ --- spike trace (SNN only), after
  exponential filter (Eq.~\ref{eq:spike-trace}).
\item $V_t \in \mathbb{R}^H$ --- membrane potential (SNN only), sampled
  from Lava \texttt{Monitor} on \texttt{LIF.v}.
\end{enumerate}

\paragraph{Per-window scalars.}
\begin{enumerate}[nosep,label=\textbf{O\arabic*.},resume]
\item Loss $\mathcal{L}$ (scalar, \texttt{float64}).
\item $\lVert z_t^{(s)} \rVert_2$ mean and std over $\mathcal{C}$
  (embedding norm).
\item $\lVert z_t^{(t)} \rVert_2$ mean and std over $\mathcal{T}$
  (teacher norm).
\item $\operatorname{Var}_{t,h}[\hat{z}_t]$ (variance across time and
  latent dimension; collapse indicator).
\item Total spike count, spike rate, synaptic event estimate.
\item Wall-clock latency components (1)--(3) from
  Section~\ref{sec:measurement}.
\end{enumerate}

\paragraph{Tolerance for round-trip comparison.}
When comparing ANN and SNN observables on identical synthetic inputs
(acceptance tests T-E1--T-E3), the following tolerances apply:
\begin{itemize}[nosep]
\item Spike trains: $\leq 1$ mismatch per 1\,000 steps per neuron.
\item Spike trace $r_t$: max absolute difference $< 10^{-4}$.
\item Readout $\hat{z}_t$: max absolute difference $< 10^{-3}$.
\item Loss: relative difference $< 10^{-6}$.
\end{itemize}

% --------------------------------------------------------------------
\subsection{Time budget estimate}\label{sec:time-budget}
% --------------------------------------------------------------------

To verify feasibility before committing compute, we decompose the
experiment into elementary units of work, measure each on the target
hardware, compose the total via closed-form equations, and attach an
uncertainty interval.  Any parameter change can be propagated
analytically without re-running benchmarks.

% ·· 5.X.1  Hardware ··············································
\paragraph{Hardware.}
Apple MacBook Pro (Model~14,6): Apple M2~Max SoC,
12~cores (8P\,+\,4E), 32\,GB unified memory,
macOS~26.1 (arm64).  All benchmarks use a single thread;
no GPU is involved (PyTorch CPU, Lava CPU simulation).

% ·· 5.X.2  Micro-benchmarks ·····································
\paragraph{Micro-benchmarks.}
Each quantity was measured after a warm-up phase (20~iterations for
PyTorch, 5~for Lava).  Timed iterations: 200 (PyTorch),
50 (Lava---graph must be rebuilt per window).  All values are
mean $\pm$ one standard deviation.

\begin{center}
\renewcommand{\arraystretch}{1.12}
\begin{tabular}{@{}llrl@{}}
\toprule
\textbf{ID} & \textbf{Quantity} & \textbf{Value} & \textbf{Unit} \\
\midrule
$t_{\text{step,ANN}}$
  & ANN training step (GRU fwd + bwd + AdamW)
  & 17.8 & ms \\
$t_{\text{step,SNN}}$
  & SNN training step (surrogate BPTT, LIF)
  & 27.4 & ms \\
$t_{\text{inf,SNN}}$
  & SNN forward-only (no grad, PyTorch LIF)
  & $11.1 \pm 0.2$ & ms \\
$t_{\text{inf,ANN}}$
  & ANN inference (GRU forward, no grad)
  & 1.8 & ms \\
$t_{\text{inf,Lava}}^{\text{run}}$
  & Lava \texttt{RunSteps(128)} (forward only)
  & $121.9 \pm 7.0$ & ms \\
$t_{\text{inf,Lava}}^{\text{build}}$
  & Lava graph construction overhead
  & $2.1 \pm 0.4$ & ms \\
$t_{\text{log}}$
  & NPZ save per window (6 tensors, $128 \times 128$)
  & 0.6 & ms \\
$t_{\text{mon}}$
  & Lava Monitor read (\texttt{get\_data} + sink)
  & $548.6 \pm 5.1$ & ms \\
\bottomrule
\end{tabular}
\end{center}

\noindent\emph{Note.}
The Lava Monitor read ($t_{\text{mon}}$) dominates per-window SNN
cost by $4.5\times$ compared to the actual \texttt{RunSteps} call.
This overhead is entirely in component~(3)
(Section~\ref{sec:measurement}) and does \emph{not} enter the primary
latency comparison.
For PyTorch training steps ($t_{\text{step,ANN}}$,
$t_{\text{step,SNN}}$), measured standard deviations were
$< 0.3$\,ms ($< 1.5$\,\%), so we report the mean only.
The SNN forward-only time $t_{\text{inf,SNN}} = 11.1$\,ms
is used for validation cost estimation
(Eq.~\ref{eq:tb-val-s}), not the full training-step time
$t_{\text{step,SNN}} = 27.4$\,ms which includes backward pass
and optimizer update.

% ·· 5.X.3  Formal time model ····································
\paragraph{Formal time model.}\label{par:time-model}
Let $S$~=~number of seeds, $P$~=~number of masking policies,
$A$~=~number of ablation conditions,
$N_{\text{tr}}$~=~training steps,
$N_{\text{eval}}$~=~validation evaluations during training
($= \lfloor N_{\text{tr}} / 500 \rfloor$),
$N_{\text{val}}$~=~validation windows,
$N_{\text{te}}$~=~test windows,
$W$~=~warm-up calls, and $N_{\text{bench}}$~=~latency benchmark calls.

For a \emph{single run} (one seed, one mask policy, one substrate):

\medskip\noindent
\textbf{ANN run:}
\begin{align}
  T_{\text{train}}^{A}
    &= N_{\text{tr}} \cdot t_{\text{step,ANN}}, \label{eq:tb-train-a} \\
  T_{\text{val}}^{A}
    &= N_{\text{eval}} \cdot N_{\text{val}} \cdot t_{\text{inf,ANN}},
       \label{eq:tb-val-a} \\
  T_{\text{teacher}}^{A}
    &= N_{\text{te}} \cdot t_{\text{inf,ANN}}, \label{eq:tb-teacher} \\
  T_{\text{test}}^{A}
    &= N_{\text{te}} \cdot
       (t_{\text{inf,ANN}} + t_{\text{log}}), \label{eq:tb-test-a} \\
  T_{\text{bench}}^{A}
    &= (W + N_{\text{bench}}) \cdot t_{\text{inf,ANN}},
       \label{eq:tb-bench-a} \\
  T_{\text{run}}^{A}
    &= T_{\text{train}}^{A} + T_{\text{val}}^{A}
     + T_{\text{teacher}}^{A} + T_{\text{test}}^{A}
     + T_{\text{bench}}^{A} + T_{\text{export}}.
       \label{eq:tb-run-a}
\end{align}

\noindent\textbf{SNN run:}
\begin{align}
  T_{\text{train}}^{S}
    &= N_{\text{tr}} \cdot t_{\text{step,SNN}}, \label{eq:tb-train-s} \\
  T_{\text{val}}^{S}
    &= N_{\text{eval}} \cdot N_{\text{val}} \cdot t_{\text{inf,SNN}},
       \label{eq:tb-val-s} \\
  T_{\text{test}}^{S}
    &= N_{\text{te}} \cdot
       \bigl(
         t_{\text{inf,Lava}}^{\text{build}}
       + t_{\text{inf,Lava}}^{\text{run}}
       + t_{\text{mon}}
       + t_{\text{log}}
       \bigr), \label{eq:tb-test-s} \\
  T_{\text{bench}}^{S}
    &= (W + N_{\text{bench}}) \cdot
       \bigl(
         t_{\text{inf,Lava}}^{\text{build}}
       + t_{\text{inf,Lava}}^{\text{run}}
       \bigr), \label{eq:tb-bench-s} \\
  T_{\text{run}}^{S}
    &= T_{\text{train}}^{S} + T_{\text{val}}^{S}
     + T_{\text{test}}^{S}
     + T_{\text{bench}}^{S} + T_{\text{export}}.
       \label{eq:tb-run-s}
\end{align}

\noindent\emph{Key distinction.}
The latency benchmark (Eq.~\ref{eq:tb-bench-s}) does \emph{not}
include $t_{\text{mon}}$: no Monitor is attached during the
1\,005-call timed loop, because the benchmark measures only
component~(1) latency.  Conversely, the test-set measurement
(Eq.~\ref{eq:tb-test-s}) \emph{does} include $t_{\text{mon}}$ and
$t_{\text{inf,Lava}}^{\text{build}}$, because every window requires a
fresh graph and full observability logging.

\medskip\noindent\textbf{Total study:}
The experimental design (Section~\ref{sec:exp-design}) defines two
tiers: a primary grid ($S \times P$ pairs) and mechanism ablations
($S \times A$ pairs, baseline mask only).
\begin{equation}\label{eq:tb-total}
  T_{\text{total}} = S \cdot P \cdot
    \bigl(T_{\text{run}}^{A} + T_{\text{run}}^{S}\bigr)
  + S \cdot A \cdot
    \bigl(T_{\text{run}}^{A} + T_{\text{run}}^{S}\bigr)
  + T_{\text{accept}}
  = S(P + A)
    \bigl(T_{\text{run}}^{A} + T_{\text{run}}^{S}\bigr)
  + T_{\text{accept}}.
\end{equation}

% ·· 5.X.4  Instantiation ········································
\paragraph{Instantiation with measured values.}
Substituting $S = 10$, $P = 3$, $A = 3$,
$N_{\text{tr}} = 5\,000$,
$N_{\text{eval}} = 10$, $N_{\text{val}} = 1\,373$,
$N_{\text{te}} = 1\,716$, $W = 5$, $N_{\text{bench}} = 1\,000$,
$T_{\text{export}} = 2.0$\,s, $T_{\text{accept}} = 300$\,s:

\begin{center}
\renewcommand{\arraystretch}{1.12}
\begin{tabular}{@{}lrrl@{}}
\toprule
\textbf{Component} & \textbf{ANN} & \textbf{SNN} & \textbf{Unit} \\
\midrule
Training ($N_{\text{tr}} = 5\,000$)
  & 89.0 & 137.0 & s \\
Validation ($10 \times 1\,373$ windows)
  & 24.7  & 152.4  & s \\
Teacher pre-compute ($1\,716$ windows)
  & 3.1  & ---   & s \\
Test-set measurement ($1\,716$ windows)
  & 4.1  & 1\,155.2  & s \\
Latency benchmark ($1\,005$ calls)
  & 1.8  & 124.6 & s \\
Export + checkpoint
  & 2.0  & 2.0   & s \\
\midrule
\textbf{Total per run} $T_{\text{run}}$
  & \textbf{124.7} & \textbf{1\,571.2} & \textbf{s} \\
  & (2.1\,min) & (26.2\,min) & \\
\midrule
\textbf{Per pair (ANN\,+\,SNN)} & \multicolumn{2}{c}{$T_{\text{pair}}
  = 1\,696.0$\,s\; (\textbf{28.3\,min})} & \\
\bottomrule
\end{tabular}
\end{center}

\noindent\emph{Dominant cost.}
SNN test-set measurement ($1\,155$\,s $= 19.3$\,min per run) accounts
for $73\,\%$ of $T_{\text{run}}^{S}$, entirely due to the per-window
Lava Monitor read ($t_{\text{mon}} = 549$\,ms).  This is a pure
\emph{observability overhead}: the Monitor captures spike traces and
membrane potentials for post-hoc analysis but does not affect the
primary latency metric (component~(1)).  The actual SNN compute
(Lava \texttt{RunSteps}) contributes only $122$\,ms/window---i.e.\
the model itself is fast, but measuring it is expensive.

\medskip
The full design (Section~\ref{sec:exp-design}) has
$S(P + A) = 10 \times (3 + 3) = 60$ pairs:
\begin{align}
  T_{\text{total}}
    &= 60 \times 1\,696.0\;\text{s}
     = 101\,760\;\text{s}
     \approx 28.3\;\text{h}. \label{eq:time-budget}
\end{align}
With acceptance tests:
$T_{\text{grand}} = 101\,760 + 300 = 102\,060$\,s
$\approx \mathbf{28.4}$\,\textbf{h}.

\paragraph{Design variants.}
If the full 28\,h budget is unacceptable, the following pre-specified
reductions are permitted (in priority order):

\begin{center}
\renewcommand{\arraystretch}{1.12}
\begin{tabular}{@{}lcccc@{}}
\toprule
\textbf{Variant} & \textbf{$S$} & \textbf{Ablations}
  & \textbf{Runs} & \textbf{Wall-clock} \\
\midrule
Full (recommended)
  & 10 & 3 ($\times$10) & 120 & 28.4\,h \\
Compromise
  & 10 & 3 ($\times$7)  & 102 & 24.1\,h \\
Minimal
  & 10 & 2 ($\times$10) & 100 & 23.6\,h \\
\bottomrule
\end{tabular}
\end{center}

\noindent The primary grid ($S=10 \times P=3 \times 2 = 60$ runs,
$\approx 14$\,h) is non-negotiable; only the ablation tier may be
reduced.  Any reduction must be documented and justified.

% ·· 5.X.5  Uncertainty propagation ······························
\paragraph{Uncertainty propagation.}
Micro-benchmark variances propagate through the linear model.
For a sum of $N$ i.i.d.\ calls with per-call variance $\sigma^2$:
$\operatorname{Var}(N \cdot X) = N^2 \sigma^2$.
The dominant stochastic components are Lava \texttt{RunSteps}
($\sigma_{\text{run}} = 7.0$\,ms) and Monitor read
($\sigma_{\text{mon}} = 5.1$\,ms).

Per-run SNN uncertainty (one seed, one policy):
\begin{align}
  \sigma_{T_{\text{test}}^{S}}
    &= N_{\text{te}}
       \sqrt{\sigma_{\text{build}}^2
           + \sigma_{\text{run}}^2
           + \sigma_{\text{mon}}^2}
     = 1716 \sqrt{0.4^2 + 7.0^2 + 5.1^2}
     \approx 14.9\;\text{s}, \label{eq:tb-sigma-test}\\
  \sigma_{T_{\text{bench}}^{S}}
    &= (W + N_{\text{bench}})
       \sqrt{\sigma_{\text{build}}^2 + \sigma_{\text{run}}^2}
     = 1005 \sqrt{0.4^2 + 7.0^2}
     \approx 7.0\;\text{s}. \label{eq:tb-sigma-bench}
\end{align}

Combined per-run SNN standard deviation (training is near-deterministic,
$\sigma < 0.3$\,ms/step):
\begin{equation}\label{eq:tb-sigma-run}
  \sigma_{T_{\text{run}}^{S}}
    = \sqrt{\sigma_{T_{\text{test}}^{S}}^2
          + \sigma_{T_{\text{bench}}^{S}}^2}
    \approx \sqrt{14.9^2 + 7.0^2}
    \approx 16.5\;\text{s}.
\end{equation}

Total study standard deviation (60~independent SNN runs):
\begin{equation}\label{eq:tb-sigma-total}
  \sigma_{T_{\text{total}}}
    = \sqrt{60} \cdot \sigma_{T_{\text{run}}^{S}}
    = 7.75 \times 16.5
    \approx 127.7\;\text{s}
    \approx 2.1\;\text{min}.
\end{equation}

\noindent\emph{Note.}
For independent runs (each seed/policy combination is a separate
process invocation), individual-run variances add, so the total
standard deviation scales as $\sqrt{n_{\text{runs}}} \cdot \sigma$,
not $n_{\text{runs}} \cdot \sigma$.

\paragraph{Pessimistic bound.}
Taking $\mu + 2\sigma$ for the total:
\begin{equation}\label{eq:tb-pessimistic}
  T_{\text{total}}^{\text{pess}}
    = 101\,760 + 2 \times 128
    = 102\,016\;\text{s}
    \approx 28.3\;\text{h}.
\end{equation}
The uncertainty is small ($< 0.3$\,\%) because
the dominant components (training, validation) are near-deterministic
and the stochastic Lava overhead, while large per-window,
averages out over 1\,716 windows.

% ·· 5.X.6  Overhead accounting ··································
\paragraph{Overhead accounting rules.}
To prevent ``hidden compute'' from invalidating the model, the
following items are explicitly assigned to components:

\begin{center}
\renewcommand{\arraystretch}{1.10}
\begin{tabular}{@{}p{4.5cm}ll@{}}
\toprule
\textbf{Activity} & \textbf{In $T_{\text{test}}$?}
  & \textbf{In $T_{\text{bench}}$?} \\
\midrule
Lava graph build ($t_{\text{build}}$)
  & yes & yes \\
Lava \texttt{RunSteps} ($t_{\text{run}}$)
  & yes & yes \\
Monitor read ($t_{\text{mon}}$)
  & yes & \textbf{no} \\
NPZ save ($t_{\text{log}}$)
  & yes & \textbf{no} \\
Loss computation
  & yes & \textbf{no} \\
Python orchestration
  & component~(3) & component~(3) \\
\bottomrule
\end{tabular}
\end{center}

\noindent The latency benchmark measures only
$t_{\text{inf,Lava}}^{\text{build}} + t_{\text{inf,Lava}}^{\text{run}}$
per call (component~(1) of Section~\ref{sec:measurement}).
The test-set measurement adds Monitor and logging because it collects
full observability data (Section~\ref{sec:observability}).

% ·· 5.X.7  Parametric sensitivity ·······························
\paragraph{Parametric sensitivity.}
The partial derivatives of $T_{\text{total}}$
(Eq.~\ref{eq:tb-total}) with respect to the main
design parameters are (at the reference point
$S = 10$, $P + A = 6$, so $n_{\text{pairs}} = 60$):

\begin{align}
  \frac{\partial T_{\text{total}}}{\partial N_{\text{tr}}}
    &= n_{\text{pairs}} \cdot
       (t_{\text{step,ANN}} + t_{\text{step,SNN}})
     = 60 \times 45.2\;\text{ms}
     = 2.71\;\text{s per extra step}, \label{eq:sens-train} \\[4pt]
  \frac{\partial T_{\text{total}}}{\partial N_{\text{te}}}
    &= n_{\text{pairs}} \cdot
       \Bigl(
         \underbrace{(t_{\text{inf,ANN}} + t_{\text{log}})}_{\text{ANN: 2.4\,ms}}
       + \underbrace{(t_{\text{build}} + t_{\text{run}}
           + t_{\text{mon}} + t_{\text{log}})}_{\text{SNN: 673.3\,ms}}
       \Bigr) \notag \\
    &= 60 \times 675.7\;\text{ms}
     = 40.5\;\text{s per extra test window}, \label{eq:sens-test} \\[4pt]
  \frac{\partial T_{\text{total}}}{\partial S}
    &= (P + A) \cdot
       (T_{\text{run}}^{A} + T_{\text{run}}^{S})
     = 6 \times 1\,696.0
     = 10\,176\;\text{s} = 2.8\;\text{h per extra seed}.
     \label{eq:sens-seed}
\end{align}

\noindent\emph{Practical consequences:}
\begin{itemize}[nosep]
\item Doubling $N_{\text{tr}}$ from 5\,000 to 10\,000 adds
  $5\,000 \times 2.71 = 13\,550$\,s ($\approx 3.8$\,h),
  total $\rightarrow 32.1$\,h.
\item Halving $N_{\text{te}}$ to 858 windows (random 50\,\% subsample)
  saves $858 \times 40.5 = 34\,749$\,s ($\approx 9.7$\,h),
  but invalidates full-test-set evaluation.
\item Adding one seed ($S = 11$) adds $2.8$\,h.
\item Parallelizing across seeds on separate P-cores (feasible on the
  8P-core M2~Max) yields near-linear speedup for training;
  with 4-way parallelism, wall-clock drops to $\approx 7$\,h.
\end{itemize}

% ·· 5.X.8  Summary box ··········································
\paragraph{Summary.}
\begin{center}
\fbox{\parbox{0.88\linewidth}{%
\textbf{Full design} ($N_{\text{tr}}\!=\!5\,000$,
$N_{\text{te}}\!=\!1\,716$, $N_{\text{bench}}\!=\!1\,000$,
$10$ seeds $\times$ $6$ conditions):\par\medskip
\centering
$T_{\text{grand}} \approx \mathbf{28\;\text{h}}$\quad
(pessimistic $\mu + 2\sigma$: $28.3$\,h)\par
\medskip
\raggedright
Primary grid only ($10 \times 3 \times 2 = 60$ runs):
$\approx 14$\,h.\par
With 4-way seed parallelism: wall-clock $\approx 7$\,h.\par
\medskip
\emph{The dominant cost is observability, not compute:
SNN test-set Monitor read accounts for 73\,\% of
$T_{\text{run}}^{S}$; the actual Lava \texttt{RunSteps}
contributes only 122\,ms/window.}
}}
\end{center}

\paragraph{Methodology.}
The micro-benchmark scripts (\texttt{benchmark\_time\_budget.py},
\texttt{benchmark\_lava.py}) are version-controlled alongside the
experiment code.  The numbers above were measured on 2025-02-07; they
will be re-measured at freeze point \texttt{freeze-v1}
(Section~\ref{sec:config}) and any discrepancy $> 10$\,\% will be
noted in the final report.

% --------------------------------------------------------------------
\subsection{Equivalence classes}\label{sec:equivalence}
% --------------------------------------------------------------------

The ANN and SNN branches must realize the \emph{same} algorithmic
function; only the \emph{substrate} (graded activations vs.\ spikes) may
differ.  The following rules define ``equivalent'' vs.\ ``deviation
(ablation)'' for every degree of freedom.

\paragraph{Positional token.}
\begin{itemize}[nosep]
\item \textbf{Equivalent:} Both branches receive the sinusoidal vector
  $\pos(t) \in \mathbb{R}^H$ (Eq.~\ref{eq:pos-enc}) with
  \emph{identical numerical values}.  Delivery mechanism may differ:
  ANN adds $\pos(t)$ to the token embedding; SNN injects $\pos(t)$ as a
  bias current into a dedicated dense process (Section~\ref{sec:lava-graph}).
\item \textbf{Deviation (= ablation):} Poisson-encoding of $\pos(t)$,
  learned positional embedding, or omitting $\pos(t)$ entirely.
\end{itemize}

\paragraph{Mask token.}
\begin{itemize}[nosep]
\item \textbf{Equivalent:} ANN replaces masked positions with a
  learned \texttt{[MASK]} embedding.  SNN replaces masked positions
  with a deterministic constant bias current of magnitude~$c_{\mask}$
  (Section~\ref{sec:mask-params}).  Both convey ``this position is
  to-be-predicted'' without leaking target information.
\item \textbf{Deviation:} Zero-masking (no signal at all), Poisson
  noise as mask, or dropout-style stochastic masking.
\end{itemize}

\paragraph{Context-only constraint (student sees only $\mathcal{C}$).}
\begin{itemize}[nosep]
\item \textbf{Equivalent:} ANN zeros out target-position tokens before
  feeding the encoder.  SNN gates input current to zero for
  $t \in \mathcal{T}$ via a deterministic binary mask on the input
  process.  Both methods must produce identical zeros in the student
  input at target positions.
\item \textbf{Deviation:} Soft masking (e.g.\ scale by 0.1 instead of
  0), or any mechanism that allows partial target information to leak.
\end{itemize}

\paragraph{General rule.}
Any modification that changes the \emph{information content} available to
the student encoder, predictor, or teacher encoder---beyond what the
substrate dictates---constitutes a deviation.  Deviations are \emph{not
forbidden}; they must be reported as separate ablation experiments with
their own identifier, seeds, and result rows.

% --------------------------------------------------------------------
\subsection{Reproducibility checklist}\label{sec:repro}
% --------------------------------------------------------------------

\begin{enumerate}[label=\textbf{R\arabic*.},nosep]
\item \textbf{Seed policy.}
  Three independent seeds are fixed and logged:
  (a)~data-split seed,
  (b)~parameter-initialization seed,
  (c)~training-order seed.
\item \textbf{Environment logging.}
  Every run records: config hash (SHA-256), git commit, Python version,
  PyTorch version, Lava version, OS, CPU model, RAM.
\item \textbf{Checkpoint policy.}
  Save student encoder, teacher encoder, and predictor at
  (a)~best validation loss,
  (b)~final step.
\item \textbf{Data versioning.}
  Dataset files are identified by SHA-256 hash; preprocessing code is
  version-controlled.
\item \textbf{Result reporting.}
  All metrics are reported as mean $\pm$ std over $S = 10$ seeds
  with 95\,\% confidence intervals and Cohen's~$d$ effect sizes
  (Section~\ref{sec:exp-design}).
\item \textbf{Determinism budget (zero-stochasticity rule).}
  The baseline experiment is \emph{fully deterministic}: no Poisson
  spike encoding, no dropout, no injected noise, no stochastic
  rounding.  Given identical seeds and inputs, two runs must produce
  bit-identical loss trajectories (tolerance: relative difference
  $< 10^{-6}$).  Any stochastic variant---Poisson encoding of inputs,
  Poisson mask token, dropout in the predictor, noise-augmented
  membrane dynamics---constitutes a \emph{separate ablation} with its
  own experiment identifier, own seeds, and its own result rows.
  Stochastic ablations additionally report the extra seed(s) they
  consume and the source of randomness.
\end{enumerate}

% --------------------------------------------------------------------
\subsection{Configuration layer and freeze point}\label{sec:config}
% --------------------------------------------------------------------

All degrees of freedom are consolidated into a single YAML configuration
file (\texttt{config/experiment.yaml}) that serves as the sole source of
truth for both implementations.  The file contains:

\begin{itemize}[nosep]
\item \texttt{data}: dataset name, $T$, $D$, sampling rate, normalization
  method, split ratios, split seed.
\item \texttt{model}: $H$, encoder type, predictor topology, $\tau$,
  $v_{\text{th}}$, $\tau_{\text{mem}}$, surrogate $k$,
  $c_{\text{mask}}$, $\alpha$ (trace decay).
\item \texttt{training}: $N_{\text{train}}$, $\text{lr}_{\max}$, warm-up
  fraction, optimizer, weight decay, masking policy + parameters.
\item \texttt{measurement}: warm-up calls, timed repetitions, thread
  count, core pinning.
\item \texttt{seeds}: data-split, init, training-order.
\end{itemize}

\paragraph{Freeze point.}
A \textbf{freeze point} is a tagged git commit
(\texttt{git tag freeze-v1}) after which no configuration value may
change.  All results reported in the final evaluation must be produced
from commits \emph{at or after} the freeze tag.  Pre-freeze runs are
labeled ``development'' and excluded from reported metrics.

% --------------------------------------------------------------------
\subsection{Acceptance tests and traceability}\label{sec:acceptance}
% --------------------------------------------------------------------

Each invariant (I1--I4), reproducibility rule (R1--R5), and failure-mode
gate (F1--F5) is linked to an automated test.  The test suite runs as a
CI pre-condition before any measurement run.

\paragraph{Invariant tests.}
\begin{enumerate}[nosep,label=\textbf{T-I\arabic*.}]
\item \textbf{Same mask.}
  For 100 random window indices, compute SHA-256 of the generated mask
  vector~$m$.  Assert ANN and SNN pipelines produce identical hashes.
\item \textbf{Same loss.}
  On a synthetic batch ($\hat{z}, z^{(t)}$ drawn from $\mathcal{N}(0,1)$),
  assert that the ANN loss function and the SNN loss function return
  values within $10^{-6}$ relative tolerance.
\item \textbf{Same inference regime.}
  Assert $B=1$ and sequential iteration in both measurement harnesses;
  log and compare window ordering.
\item \textbf{Same latency definition.}
  Assert that both harnesses use the same timing wrapper
  (\texttt{shared/timing.py}) and report identical column names.
\end{enumerate}

\paragraph{Export tests.}
\begin{enumerate}[nosep,label=\textbf{T-E\arabic*.}]
\item \textbf{Spike-train round-trip.}
  Run 1\,000 steps on a random input in PyTorch LIF and Lava LIF with
  exported weights.  Assert $\leq 1$ spike mismatch per neuron.
\item \textbf{Trace round-trip.}
  Compute spike trace $r_t$ (Eq.~\ref{eq:spike-trace}) from the spike
  trains of T-E1 in both frameworks.  Assert max absolute difference
  $< 10^{-4}$.
\item \textbf{Readout round-trip.}
  Compute $\hat{z}_t$ (Eq.~\ref{eq:readout}) from both traces.  Assert
  max absolute difference $< 10^{-3}$.
\end{enumerate}

\paragraph{Failure-mode gates.}
\begin{enumerate}[nosep,label=\textbf{T-F\arabic*.}]
\item \textbf{Collapse gate.}
  After each epoch, if $\operatorname{Var}[\hat{z}_t] < 10^{-3}$ for
  3~consecutive epochs, the run is flagged \texttt{INVALID} and halted.
  Mitigation runs (with adjusted $\tau$ or regularizer) are logged
  separately and \emph{never} mixed into primary results.
\item \textbf{Saturation gate.}
  If mean spike rate $> 0.9$ over an epoch, flag and halt.
\item \textbf{Silence gate.}
  If total spike count $= 0$ for any test window, flag and halt.
\end{enumerate}

\paragraph{Traceability matrix.}
Table~\ref{tab:traceability} maps each specification item to its
implementing module, test, and log artifact.

\begin{table}[htbp]
\centering
\caption{Traceability matrix: specification $\to$ module $\to$ test
$\to$ artifact.}
\label{tab:traceability}
\smallskip
\renewcommand{\arraystretch}{1.1}
\begin{tabular}{@{}llll@{}}
\toprule
\textbf{Spec} & \textbf{Module} & \textbf{Test} & \textbf{Artifact} \\
\midrule
I1 (same task)
  & \texttt{shared/data.py}
  & T-I1
  & mask hash log \\
I2 (same target)
  & \texttt{shared/loss.py}
  & T-I2
  & loss comparison log \\
I3 (same regime)
  & \texttt{shared/harness.py}
  & T-I3
  & window-order log \\
I4 (same latency)
  & \texttt{shared/timing.py}
  & T-I4
  & timing CSV \\
R1 (seeds)
  & \texttt{config/experiment.yaml}
  & CI seed check
  & config hash \\
R2 (env logging)
  & \texttt{shared/env\_log.py}
  & CI env check
  & \texttt{env.json} \\
R3 (checkpoints)
  & training scripts
  & file-exists assert
  & \texttt{*.pt} files \\
R4 (data version)
  & \texttt{data/download.py}
  & SHA-256 assert
  & hash log \\
F1 (collapse)
  & \texttt{shared/gates.py}
  & T-F1
  & variance log \\
F2 (saturation)
  & \texttt{shared/gates.py}
  & T-F2
  & spike-rate log \\
F3 (silence)
  & \texttt{shared/gates.py}
  & T-F3
  & spike-count log \\
\bottomrule
\end{tabular}
\end{table}

% --------------------------------------------------------------------
\subsection{Known failure modes and mitigations}\label{sec:failure}
% --------------------------------------------------------------------

\begin{enumerate}[label=\textbf{F\arabic*.},nosep]
\item \textbf{Representation collapse.}
  Loss decreases but latent variance $\operatorname{Var}[\hat{z}_t]
  \to 0$.\\
  \emph{Mitigation:} monitor latent variance per epoch; if below
  threshold, add variance regularizer
  $\lambda\!\cdot\!(1/\operatorname{Var}[\hat{z}_t])$ or reduce
  $\tau$.

\item \textbf{SNN spike saturation.}
  Spike rate $\to 1$ (every neuron fires every step).\\
  \emph{Mitigation:} increase firing threshold $v_{\text{th}}$;
  reduce input scaling; verify membrane leak $\mathrm{d}u > 0$.

\item \textbf{SNN silence (zero activity).}
  No spikes emitted; readout is constant.\\
  \emph{Mitigation:} increase input scaling; decrease
  $v_{\text{th}}$; check weight initialization variance.

\item \textbf{ANN trivial prediction.}
  Model outputs a constant vector (mean of targets).\\
  \emph{Mitigation:} check embedding norm distribution and
  per-timestep variance; verify masking is applied correctly.

\item \textbf{Lava runtime hang (macOS).}
  Bit-accurate Loihi model deadlocks with \texttt{fork}-based
  multiprocessing on Apple Silicon.\\
  \emph{Mitigation:} use
  \texttt{Loihi2SimCfg(select\_tag="floating\_pt")} and add a
  safety timeout.
\end{enumerate}

% ====================================================================
\section{Outcome Interpretation Framework}\label{sec:outcomes}
% ====================================================================

To prevent post-hoc ``storytelling,'' we pre-register the complete set of
possible experimental outcomes and the exact conditions under which each
scientific conclusion is warranted.  All conditions are evaluated
\emph{after} applying the Holm--Bonferroni correction
(Section~\ref{sec:exp-design}, S4).

\paragraph{Notation.}
For a given masking policy~$p$ and seed~$s$, denote:
\begin{itemize}[nosep]
\item $\mathcal{L}_{s,p}^{X}$~---~prediction loss
  (Eq.~\ref{eq:loss}) on the full test set, substrate
  $X \in \{\text{ANN}, \text{SNN}\}$;
\item $N_{\text{spikes},s,p}$,
  $N_{\text{syn},s,p}$~---~spike count and synaptic events
  (SNN only);
\item $\hat{E}_{s,p}$~---~energy proxy
  (Eq.~\ref{eq:energy});
\item $\operatorname{Var}[\hat{z}]_{s,p}$~---~latent variance;
\item $\tilde{t}_{s,p}^{X}$~---~median wall-clock latency.
\end{itemize}
All metrics are aggregated as mean~$\pm$~std over $S = 10$ seeds
and tested via two-sided paired $t$-test at $\alpha = 0.05$
(with Wilcoxon robustness check).

% ·· 6.1  Scenario table ·········································

\begin{table}[htbp]
\centering
\caption{Pre-registered outcome scenarios.  Each row defines a
  mutually exclusive interpretation; exactly one will be reported.
  $p_{\mathcal{L}}$ denotes the adjusted $p$-value for the paired
  loss comparison;
  $r_{\text{spike}} = N_{\text{spikes}} / (H \times T)$ is the
  mean spike rate.}
\label{tab:outcomes}
\smallskip
\renewcommand{\arraystretch}{1.15}
\begin{tabular}{@{}cp{3.5cm}p{6.5cm}@{}}
\toprule
\textbf{ID} & \textbf{Conditions} & \textbf{Conclusion} \\
\midrule
\textbf{A}
  & $p_{\mathcal{L}} > 0.05$
    \emph{and} $\hat{E}_{\text{SNN}} \ll
    \text{MACs}_{\text{ANN}}$
  & \emph{Neuromorphic equivalence with higher efficiency.}
    JEPA objective is realizable at comparable quality via sparse
    event-driven dynamics; JEPA is a neuromorphic-native workload. \\[4pt]
\textbf{B}
  & $p_{\mathcal{L}} < 0.05$,
    $\bar{\mathcal{L}}^{\text{SNN}} >
     \bar{\mathcal{L}}^{\text{ANN}}$
    \emph{and}
    $\hat{E}_{\text{SNN}} \ll
    \text{MACs}_{\text{ANN}}$
  & \emph{Accuracy--efficiency trade-off.}
    SNN approximates JEPA with significantly lower computational
    activity; classic quality-vs-cost trade-off. \\[4pt]
\textbf{C}
  & $p_{\mathcal{L}} < 0.05$,
    $\bar{\mathcal{L}}^{\text{SNN}} <
     \bar{\mathcal{L}}^{\text{ANN}}$
  & \emph{SNN superiority.}
    Temporal JEPA is more naturally suited to dynamic neural
    representations than to tensor-sequential processing. \\[4pt]
\textbf{D}
  & $p_{\mathcal{L}} > 0.05$
    \emph{and} $r_{\text{spike}} > 0.8$
  & \emph{Loss of sparsity.}
    JEPA objective forces near-continuous activity, negating
    the neuromorphic advantage; JEPA is \emph{not} a
    neuromorphic-native workload. \\[4pt]
\textbf{E}
  & $\operatorname{Var}[\hat{z}] < 10^{-3}$
    \emph{or} $r_{\text{spike}} > 0.9$
    \emph{or} total spike count $= 0$
  & \emph{Training failure.}
    Surrogate-gradient BPTT is insufficient for JEPA on the SNN
    substrate under the tested hyperparameters.
    Run is marked \texttt{INVALID} per failure gates F1--F3
    (Section~\ref{sec:failure}). \\[4pt]
\textbf{F}
  & $p_{\mathcal{L}} > 0.05$
    \emph{and}
    $\hat{E}_{\text{SNN}} \approx
    \text{MACs}_{\text{ANN}}$
  & \emph{Task insufficiency.}
    The UCI~HAR benchmark is not computationally demanding enough
    to differentiate the two paradigms. \\
\bottomrule
\end{tabular}
\end{table}

\paragraph{Disambiguation rules.}
The scenarios are evaluated in order A--F.
Scenario~E takes priority: if any run triggers a failure gate,
it is excluded and the remaining seeds are tested for
A--D/F (minimum $S_{\text{eff}} \geq 5$ valid seeds required;
otherwise the experiment is declared \texttt{INCONCLUSIVE}).
Between A and D, the discriminator is the spike rate
$r_{\text{spike}}$: if $r_{\text{spike}} \leq 0.8$, scenario~A
applies; otherwise scenario~D.
Between A and F, the discriminator is the energy ratio
$\hat{E}_{\text{SNN}} / \text{MACs}_{\text{ANN}}$: if
$< 0.5$, scenario~A; if $\in [0.5, 1.5]$, scenario~F.

\paragraph{Consistency requirement.}
A conclusion is reported only if the same scenario holds for
\emph{all three} masking policies.  If policies yield different
scenarios (e.g.\ A for future-block but B for random-drop), we
report the per-policy results and classify the overall outcome as
\textbf{mixed}, with the dominant scenario (by majority)
highlighted.

\paragraph{Effect-size thresholds.}
For energy-ratio claims (``$\ll$''), we require
$\hat{E}_{\text{SNN}} / \text{MACs}_{\text{ANN}} < 0.5$
(at least $2\times$ reduction).  For spike-rate claims
(``high'' vs.\ ``sparse''), we use
$r_{\text{spike}} \leq 0.3$ as ``sparse'' and
$r_{\text{spike}} > 0.8$ as ``dense.''
The intermediate range $[0.3, 0.8]$ is reported as
``moderate sparsity'' and weakens the A/D conclusion
to ``suggestive.''

\paragraph{Falsifiability.}
Each scenario is falsifiable by construction:
scenario~A is falsified if $p_{\mathcal{L}} < 0.05$ or
$\hat{E}$ ratio $\geq 0.5$;
scenario~C is falsified if $\bar{\mathcal{L}}^{\text{SNN}}
\geq \bar{\mathcal{L}}^{\text{ANN}}$;
scenario~E is falsified by all runs passing failure gates.
No outcome can be ``spun'' into a different narrative without
violating the pre-registered conditions above.

\section{Conclusion}

We have formalized a minimal experimental framework for comparing a
JEPA-style predictive objective implemented on two fundamentally different
computation paradigms: a classical ANN (GRU) and a spiking neural network
simulated in Intel Lava.  Both implementations share the same input data,
the same latent-prediction objective~\eqref{eq:loss}, and the same
evaluation metrics, enabling a fair comparison of prediction quality,
latency, and computational effort.  The implementation
specification (Section~\ref{sec:impl-spec}) locks down the task
contract, baseline configurations, training protocol, and measurement
procedures to ensure that experimental results are reproducible and free
from accidental asymmetries between the two paradigms.
The pre-registered outcome interpretation framework
(Section~\ref{sec:outcomes}) guarantees that every possible experimental
result maps to a well-defined scientific conclusion, making the study
falsifiable by design.

The key research question remains: \emph{Is the JEPA-style predictive
world-model objective more naturally realizable using spiking neural
networks?}  The experimental setup described in this document---with
$S = 10$ seeds, three masking policies, three mechanism ablations, and
a pre-specified statistical plan---provides the tools to answer it.

\bibliographystyle{plainnat}
\bibliography{references}

\end{document}
